\scanpage{55}
\mspage{10}
% archive sheet number: 54
\deleted{\unsure{utilise donc le|utilise alors le|utilisé dans la}}

\deleted{\uline{Lemme} Soit $X$ fini sur un corps $k$, $k'$ une extension primaire de $k$, $\unsure{H'|X'}$ un sous-schém. (nécessairement \unsure{fermé|fini}) de $X'=X\otimes_k k'$. Alors l'application \unsure{est canonique}}

Comme $k'/k$ est primaire, \inserted{et $G/k$ fini}, l'application
$\deleted{\unsure{X'}}G'\longrightarrow G$ est bijective (mettre un
\uline{lemme}) donc en tant qu'ensemble
$H'$ est l'image inverse d'un ss-ensemble
\deleted{\unsure{de}} $H$ de $G$. \deleted{\unsure{qui est muni de la}}
\inserted{D'ailleurs}, \deleted{\unsure{puisque que}} par
\deleted{\unsure{de}} une extension radicielle finie convenable
$k_1$ de $k$, \deleted{\unsure{on trouve une} fois que}
\inserted{($G\otimes_k k'$, \unsure{d'id|d'id.} \unsure{sep.} séparable sur $k'$ (\unsure{sic|vif})) donc}
\deleted{que} $G_1\unsure{:|=}G\otimes_k k_1$ a des extensions résiduelles
séparables. \deleted{En} Si donc $H_1$ est l'image
inverse de $H$ dans $G_1$, muni de
la structure \uline{réduite} induite, $H_1$ est
un \uline{sous-schéma} \inserted{\uline{en groupes}} fermé de $G_1$.

De plus, $H_1\otimes_{\unsure{k'_1|k_1}}k'_1$ a $=$ espace
sous-jacent que $H'_1=H'\otimes_k k'_1$. On a
fini.
